% /root/monodromy, 2026-11-19. Proved; independent root and reduction_descent audits passed.
\subsection{A supersingular product: a double Hasse component and a full image}
\label{subsec:supersingular-product-degree-one}
\status{Proved and checked. The parent and reduction/descent agent independently
audited the complete determinant, local wild-inertia and full-group proofs.
Both independently read and reran the polynomial/Pfaffian certificate and
the separate finite-field group verifier; all checks passed.}

Let
\[
 E:\ y^2+y=x^3\quad\text{over }\mathbf F_2,
 \qquad X=E\times E,\qquad
 L=\mathcal O_E(3O)\boxtimes\mathcal O_E(3O).
\]
The indicated Segre embedding identifies ambient hyperplane sections with
the complete system $P=\mathbf P(W)$, where $W=H^0(X,L)$ has dimension nine.
Write $s=(s_{ij})_{0\le i,j\le2}$, in lexicographic order, in the tensor
basis of $(1,x,y)$, and write $S=(s_{ij})$ for the corresponding
$3\times3$ matrix. Put $J_3=(\delta_{i,2-j})$ and $J=J_3\otimes J_3$.

\begin{proposition}[An explicit nonreduced relative Hasse divisor]
\label{prop:ss-product-relative-Hasse}
The relative coherent Frobenius bundle is $\mathcal E=\Omega^1_P$.
Its determinant section, up to the constant trivialization implicit in
the chosen basis, is
\begin{equation}\label{eq:ss-product-Hasse-factorization}
 h(s)=D(s)^2Q(s),\qquad
 D(s)=\det S,\qquad Q(s)=s^{[2]t}Js.
\end{equation}
The cubic $D$ is geometrically irreducible; the cubic $Q$ is
geometrically smooth and irreducible; they are distinct. Thus the
relative Hasse divisor has a component of multiplicity two and a
different component of multiplicity one. The generic relative stable
rank is eight. A smooth generic section has genus ten and $2$-rank eight.
\end{proposition}

\begin{proof}
Use $\eta=dx$ on $E$. Writing a rational function as $a^2+xb^2$
gives $C((a^2+xb^2)\,dx)=b\,dx$. The curve equation gives
\[
 C(1)=0,\quad C(x)=1,\quad C(y)=x,\quad
 C(x^2)=C(y^2)=0,\quad C(xy)=y.
\]
Consequently the root-removed Cartier multiplication matrix on the
degree-three one-factor space is
\[
 M_E(u)=\begin{pmatrix}u_1&u_0&0\\u_2&0&u_0\\0&u_2&u_1\end{pmatrix},
 \qquad
 J_3M_E(u)=\begin{pmatrix}0&u_2&u_1\\u_2&0&u_0\\u_1&u_0&0\end{pmatrix}.
\]
If $M_i$ is the coefficient of $u_i$, then the full tensor matrix is
$M(s)=\sum s_{ij}M_i\otimes M_j$. Direct differentiation of the
six-term determinant shows
\begin{equation}\label{eq:ss-Hessian}
 A(s):=JM(s)=\operatorname{Hess}_s(D),\qquad
 A(s)^t=A(s),\quad A(s)_{ii}=0,\quad M(s)s=0.
\end{equation}
The eight-dimensional Pfaffians of the alternating $9\times9$
matrix $A(s)$ form the vector
\[
       \operatorname{Pf}_8(A(s))=D(s)s.
\]
This identity may be verified by expansion of the nine Pfaffians.
The retained exact polynomial program checks all nine over
$\mathbf F_2[s_0,\ldots,s_8]$, without specialization. The standard
Pfaffian adjugate identity, or the same expansion of cofactors, gives
\[
 \operatorname{adj}M(s)=D(s)^2s\,(Js)^t,
 \qquad
 \operatorname{adj}(M(s)^t)=D(s)^2(Js)s^t.
\]

Kunneth gives $H^1(X,L^{-1})=0$, and the hypersurface exact sequence
gives the Euler sequence
\[
 0\longrightarrow\mathcal E\longrightarrow
 W^*\otimes\mathcal O_P(-1)\longrightarrow\mathcal O_P\longrightarrow0.
\]
The quotient Frobenius is zero because $X$ is supersingular.
In affine coefficient coordinates the relative Frobenius is therefore
the restriction of $M(s)^t$ from $(s^{[2]})^\perp$ to $s^\perp$.
The cofactor formula identifies the determinant of this restriction
with $D(s)^2(s^{[2]})^tJs$: indeed the exterior-product vector of a
basis of the source hyperplane is its normal $s^{[2]}$, and the
cofactor matrix carries it to
$D(s)^2s\,(Js)^ts^{[2]}$, a scalar multiple of the target normal $s$.
The projective twists agree: $\det\mathcal E=\mathcal O_P(-9)$, so
the Frobenius determinant is a section of $\mathcal O_P(9)$.
The displayed homogeneous polynomial has degree nine, proving
\eqref{eq:ss-product-Hasse-factorization} globally.

The rank-at-most-two matrix locus is the image closure of the
irreducible space of products of a $3\times2$ and a $2\times3$
matrix, and is the determinant hypersurface. Thus $D$ is
geometrically irreducible. Moreover
\[
                  \nabla Q(s)=Js^{[2]},
\]
which never vanishes at a projective point. Thus $Q$ is smooth;
two distinct positive-degree projective hypersurfaces in $\mathbf P^8$
intersect, so smoothness also implies geometric irreducibility.
At $S=\operatorname{diag}(0,1,1)$ one has $D=0,Q=1$, proving that
the two cubics are distinct. At $S=I_3$ both $D$ and $Q$ are one.
Thus the relative Frobenius is generically invertible.
The smooth open is nonempty by the very ample embedding, so it meets
this invertibility open. Adjunction gives $g=1+L^2/2=10$.
The ambient $H^1(X,\mathcal O_X)$ has nilpotent Frobenius and the
relative stable rank is eight, so the smooth generic curve has
$2$-rank eight.
\end{proof}

\begin{lemma}[A first-derivative criterion for wild transvection inertia]
\label{lem:ss-simple-pole-transvection}
Let $R=k[[t]]$ with $k$ algebraically closed of characteristic two,
and let a rank-$r$ vector bundle over $R$ have regular semilinear
Frobenius $F$. Suppose that its special fibre has stable rank $r-1$,
its linearized kernel has dimension one, and its determinant has
valuation one. Let $u$ generate the semilinear kernel in the special
fibre, and let $v$ be the first derivative of $F(\widetilde u)$ at
$t=0$, for a regular lift $\widetilde u$. If $v\notin ku$, the
generic Lang local system has a nontrivial transvection in its inertia.
\end{lemma}

\begin{proof}
The finite part of the Lang scheme supplies $r-1$ regular fixed
sections, whose special values span the stable subspace. This is the
finite-part lemma already proved using Stacks~\cite{Stacks0BTN,Stacks09ZS}.
They form a Frobenius-fixed direct summand. Complete them by a lift of
$u$. In this frame
\[
 F=\begin{pmatrix}I_{r-1}&b(t)\\0&a(t)\end{pmatrix},
 \qquad a(t)=a_1t+O(t^2),\qquad b(t)=b_1t+O(t^2),\qquad a_1\ne0.
\]
The constant column is zero because $u$ is in the special kernel.
Changing the lift changes its Frobenius image only in order at least
two, so its first derivative is well defined in the special fibre.
The hypothesis $v\notin ku$ is precisely $b_1\ne0$.

The nonzero fixed vector in the one-dimensional generic quotient has
coordinate $z=a(t)^{-1}$. Lifting it requires solving
\[
                  y_i^2+y_i=b_i(t)/a(t)^2.
\]
At least one right side has a nonzero simple pole. Such an
Artin--Schreier equation is ramified: if $y$ lay in an unramified
extension, its negative valuation would give an even pole order
for $y^2+y$, a contradiction; a nonnegative valuation gives no pole.
Over the strictly henselian field all inertia fixes the regular
$(r-1)$-dimensional Lang subspace and its one-dimensional quotient.
A nonidentity element therefore differs from the identity by a rank-one
square-zero map, which is a transvection.
\end{proof}

\begin{theorem}[Full degree-one image on a supersingular product]
\label{thm:ss-product-degree-one-full-group}
For $X,L$ above, over every perfect field of characteristic two,
the degree-one relative Hasse--Witt local system on its maximal-rank
smooth open has
\[
                   G_{\rm geom}=G_{\rm arith}=\GL_8(\mathbf F_2).
\]
In particular its geometric representation is irreducible, despite
the double component of its coherent Hasse divisor.
\end{theorem}

\begin{proof}
We give two small exact certificates over
$k_8=\mathbf F_2[\alpha]/(\alpha^3+\alpha+1)$.
For readability an integer $c_0+2c_1+4c_2$ denotes the field element
$c_0+c_1\alpha+c_2\alpha^2$; all coefficient lists below use this
encoding and the fixed lexicographic tensor basis.

First take
\[
                s=(6,7,0,5,4,1,6,0,4).
\]
The relative Frobenius is invertible and the characteristic polynomial
of its third iterate, up to replacing it by its reciprocal according
to the chosen Frobenius convention, is
\begin{equation}\label{eq:ss-degree-one-Frobenius-polynomial}
                    T^8+T^7+T^5+T+1.
\end{equation}
The retained verifier checks its irreducibility over $\mathbf F_2$
by the Frobenius-power and gcd criterion and checks that it is not
reciprocal. Thus the arithmetic image contains the element required
in Lemma~\ref{lem:product-frobenius-group-test}.
The verifier uses an independent restriction-of-scalars calculation:
if $B=A A^{[2]}A^{[4]}$ is the full third Frobenius iterate, then
$B^{[2]}=A^{[2]}A^{[4]}A$ has the same characteristic polynomial,
because cyclic products of square matrices do. Thus that polynomial
has coefficients in $\mathbf F_2$. Restriction from $\mathbf F_8$ to
$\mathbf F_2$ cubes it. The directly computed $27\times27$ integer
characteristic polynomial, reduced modulo two, is
$(T(T^8+T^7+T^5+T+1))^3$; unique factorization, with monic factors,
therefore proves the asserted nine-dimensional polynomial and its
eight-dimensional relative factor.

For a geometric transvection, take
\[
 s_0=(0,6,4,5,5,5,3,3,4),\qquad\delta=(1,0,\ldots,0),
 \qquad s(t)=s_0+t\delta.
\]
Here $D(s_0)\ne0$, $Q(s_0)=0$, and $dQ_{s_0}(\delta)=6\ne0$.
The full linearized matrix $M(s_0)^t$ has rank eight; its semilinear
stable rank is seven. A generator of its semilinear kernel is
\[
 u=(2,7,7,3,3,3,2,4,0)=(Js_0)^{[1/2]}.
\]
The verifier checks $s_0^tu=0$, so this is also the relative kernel.
The derivative on this kernel along the displayed affine line is
\[
 v=M(\delta)^tu^{[2]}=(0,0,0,0,4,3,0,5,5).
\]
These two vectors are independent. The determinant factorization
shows that the relative determinant has valuation one along this
line at $t=0$.

For completeness, this ambient calculation is also the derivative
in a regular frame of the relative bundle. A lift of $u$ in
$s(t)^\perp$ has square $u^{[2]}+O(t^2)$; hence its image derivative
is exactly $M(\delta)^tu^{[2]}$. A change of output frame has no
first-order effect because its image at $t=0$ is zero. Passing
from the affine coefficient cone to a projective chart only subtracts
a radial tangent vector; that adds a multiple of
$M(s_0)^tu^{[2]}=0$ and changes neither derivative test.
Lemma~\ref{lem:ss-simple-pole-transvection} therefore supplies a
nontrivial geometric transvection.

Let $H$ be its normal closure in the arithmetic image $G$.
Then $H$ is nontrivial, generated by transvections, normal in $G$,
and lies in the geometric image. Applying
Lemma~\ref{lem:product-frobenius-group-test} with $N=8$ proves
$H=G=\GL_8(\mathbf F_2)$. That lemma explicitly handles Clifford
blocks using the primitive prime $17$ and excludes all form-preserving
alternatives by nonreciprocity; its only classification input is
\cite[Theorem~3.6, printed p.~22]{KantorSeress2001}.

All calculations were made on the invertibility open of the coherent
bundle over the full projective parameter space. Restriction to the
nonempty smooth open does not change the image: a connected finite
etale cover of a normal integral scheme remains connected on a
dense open. Thus the use of coefficient points or a local line does
not require a separate smoothness assertion for those particular
sections. Finally geometric connectedness of the finite frame cover
is preserved by algebraically closed base extension. The same
geometric full group, and hence the same arithmetic full group, holds
over every perfect base field of characteristic two.
\end{proof}

\begin{remark}[Retained exact calculations]
The polynomial certificate is
\path{agents/monodromy/code/supersingular_product_degree_one.py}, with
output \path{agents/monodromy/outputs/supersingular_product_degree_one.json}.
It verifies all nine Pfaffian identities and all $511$ projective
$\mathbf F_2$ coefficient vectors. The extension-field discovery
calculation is
\path{agents/monodromy/code/supersingular_product_frobenius_probe.py};
its output retains the field polynomials, coefficients, characteristic
polynomials and local derivative data. A separate certificate verifier
\path{agents/monodromy/code/verify_supersingular_degree_one_group.py}
checks only the two stated $\mathbf F_8$ points and does not depend
on the discovery search. Its output is
\path{agents/monodromy/outputs/verify_supersingular_degree_one_group.json}.
\end{remark}
